% Sección · Mandato acordado (de la E1, condensado en la E2 para respetar las 10 páginas).
% Todas las cifras están verificadas contra output/cifras.md (componente B).

\section{Mandato acordado}
\label{sec:mandato}

Cada hallazgo va marcado con su origen: \textbf{(V)} verbalizado por la contraparte, \textbf{(M)} medido por nuestro análisis, \textbf{(I)} inferido por el equipo.

\paragraph{La pregunta.} De las cuatro del README, trabajamos la segunda: el precio promedio por litro que subió, ¿es porque subimos precios o porque vendimos mayor proporción de gran reserva?

\paragraph{Por qué esa.} Sale de la reunión del 4 de septiembre. Sin que se lo preguntáramos, la contraparte dijo que el encargo nace de la conversión entre cajas y litros: «nos dimos cuenta que teníamos un problema y por eso los llamamos» (V). Nombró el vacío que quiere cerrar, «no se disgrega cepa, reserva, gran reserva ni varietal al tomar la decisión» (V), y una fractura de gobierno de datos: «ellos hablan de cajas, nosotros hablamos de litros» (V), porque comercial vende cajas y operaciones produce litros. Medir por litro es decisión nuestra, no una instrucción del cliente: es la única unidad que compara precios entre formatos de caja (I). El dato confirma que la pregunta no es obvia: el precio medio por litro de exportación subió 11,2\,\% entre semestres de 2025, de 4,574 a 5,085 USD/L, mientras los precios de lista bajaron entre 0,15\,\% y 0,19\,\% en los tres tiers y la gran reserva pasó de 22,7\,\% a 32,5\,\% de los litros (M).

\paragraph{Por qué no las otras.} La primera, el tipo de cambio, queda como contexto: la contraparte no lo planteó como problema y declaró no tener política formal para valorizar exportaciones (V). La tercera, margen por mercado, no se puede responder con lo entregado: faltan aranceles, costos de despacho y el tramo de dólar a euro de los clientes europeos (V). La cuarta, con qué fecha convertir a pesos, no es otra pregunta sino una decisión de definición de la nuestra, y la tratamos en la sección~\ref{sec:metrica}.

\paragraph{El hilo de correo y la definición que quedó abierta.} El hilo se abrió el 8 de septiembre con el acta. Las consultas enviadas por esa vía no fueron respondidas, así que la equivalencia de caja a litro, el tratamiento del monto neto y la estacionalidad del alza quedan como supuestos cuantificados en la bitácora (S1, S3 y S8). La equivalencia era la definición abierta: operaciones afirmó que «no todas las cajas tienen el mismo número de litros» (V), pero el maestro de costos declara 9,0 litros para los dos formatos, 12$\times$750 y 6$\times$1500 (M), y el precio de lista por caja del magnum, USD 77,26, coincide con el del gran reserva estándar, USD 77,20 (M), lo que solo se sostiene si ambos rinden los mismos litros (I). La cerramos en 9,0 litros: cuando el cliente y el dato no coinciden, adoptamos el dato si otra fuente lo verifica, y aquí la verificación viene de otro archivo. Con 18 o 4,5 litros por caja magnum la mezcla sigue explicando el 108\,\% del alza (anexo~D).
